% ===== TITLE OPTIONS — uncomment one =====
% Option 1 (current favorite — the punchline):
\newcommand{\papertitle}{Linear Methods Fail Where They Matter Most: \\
Geometric Dissociation and Causal Subspaces in Neural Populations}
% Option 2 (provocative question):
%\newcommand{\papertitle}{Why Does Your Metric Disagree With Mine? \\
%Dimensionality Mediates Geometric Dissociation in the Brain}
% Option 3 (clean and direct):
%\newcommand{\papertitle}{The Metric Matters: Dimensionality Structure \\
%Determines Which Neural Similarity Measures Are Informative}
% Option 4 (causal framing forward):
%\newcommand{\papertitle}{Causal Subspaces Are Invisible to Linear Methods: \\
%Evidence from Brain-Wide Neuropixels Recordings}
% Option 5 (structured VAE angle):
%\newcommand{\papertitle}{Structured Disentanglement Reveals Causal Choice \\
%Pathways Hidden from Linear Discriminant Analysis}
% Option 6 (broadest — geometry + causality):
%\newcommand{\papertitle}{Neural Geometry Is Not Metric-Neutral: \\
%Dimensionality, Dissociation, and Causal Subspace Discovery}
% Option 7 (shortest):
%\newcommand{\papertitle}{Geometric Dissociation in Neural Population Codes}
% Option 8 (silencing angle):
%\newcommand{\papertitle}{Optogenetic Ground Truth Validates Nonlinear \\
%Causal Subspaces Across the Mouse Brain}
% Option 9 (the finding, stated plainly):
%\newcommand{\papertitle}{A 3.4$\times$ Causal Signal That Linear Methods Miss: \\
%Nonlinear Subspace Discovery in Neural Populations}
% Option 10 (dual-finding):
%\newcommand{\papertitle}{Two Geometric Dissociations: Metrics Disagree \\
%About Neural Similarity, and Linear Methods Underestimate Causal Structure}

\documentclass{article}

% TMLR format
\usepackage[accepted]{tmlr}

\usepackage{amsmath,amssymb,amsfonts}
\usepackage{booktabs}
\usepackage{graphicx}
\usepackage{hyperref}
\usepackage{multirow}
\usepackage{xcolor}

% Placeholder command for pending results
\newcommand{\pending}[1]{\textcolor{red!70!black}{\textbf{[PENDING: #1]}}}

\title{\papertitle}

\author{
Elliot Tower \\
\texttt{elliot@elliottower.com}
}

\begin{document}
\maketitle

\begin{abstract}
We study how the choice of geometric similarity metric interacts with brain region dimensionality, and how this interaction conceals causal structure from standard linear methods. Using Neuropixels recordings across 73 brain regions in 10 mice \citep{steinmetz2019distributed}, we establish two geometric dissociations. \textbf{First}, linear (CKA) and nonlinear (UMAP Procrustes) similarity are strongly anti-correlated ($\rho = -0.85$, 95\% CI $[-0.87, -0.83]$), and this dissociation is fully mediated by effective dimensionality. \textbf{Second}, a structured variational autoencoder discovers causal choice subspaces that are 3.4$\times$ stronger than those found by linear discriminant analysis (LDA), winning in 73/73 brain regions ($p = 5.7 \times 10^{-14}$). This advantage generalizes across four task variables (choice, evidence strength, feedback, stimulus side; all $p < 10^{-13}$) and is validated against optogenetic silencing ground truth: LDA-identified subspaces are \emph{anti-correlated} with causal importance ($\rho = -0.73$, $p = 0.01$), while VAE subspaces show a positive trend, with the method difference significant (bootstrap $\Delta\rho = +1.06$, 95\% CI $[0.31, 1.51]$). The engagement and choice subspaces are near-orthogonal ($d_G = 1.82$), confirming that the VAE disentangles distinct cognitive variables. However, we also show that the VAE's interchange intervention accuracy (IIA) metric is potentially vacuous for nonlinear methods---achieving comparable scores on random noise as on real data---motivating external validation over metric-internal evidence. These findings establish that linear subspace methods systematically fail in the highest-dimensional brain regions where causal signals are strongest, and that external causal validation (not IIA alone) is necessary to assess subspace quality.
\end{abstract}

\section{Introduction}

A growing body of work characterizes neural computation through the geometry of population activity \citep{cunningham2014dimensionality, gallego2017neural}. When $N$ neurons fire together, their joint activity traces trajectories in an $N$-dimensional space, and the low-dimensional structure of those trajectories reveals the underlying computation. Decision variables live in low-dimensional subspaces \citep{steinmetz2019distributed}. Motor plans occupy orthogonal manifolds \citep{gallego2017neural}. Representations drift over time within fixed-dimensional structures \citep{driscoll2022representational}.

This geometric view raises a methodological question with substantive consequences: \emph{does the choice of analysis method change the scientific conclusion?} The field's standard tools for identifying task-relevant subspaces---linear discriminant analysis (LDA), principal component analysis (PCA), and their variants---assume that the causal structure of neural population codes lies in linear subspaces. If this assumption fails, these methods will not merely be imprecise; they will systematically misidentify which brain regions carry the strongest causal signals.

We test this directly. Using 1,316 cross-session pairs across 73 brain regions in the Steinmetz et al.\ (2019) Neuropixels dataset, we find two geometric dissociations that together demonstrate the failure of linearity assumptions:

\textbf{Dissociation 1: Metrics disagree.} Linear (CKA) and nonlinear (UMAP Procrustes) similarity are strongly anti-correlated ($\rho = -0.85$). This is mediated by effective dimensionality: high-dimensional regions appear dissimilar to linear methods but highly similar to nonlinear ones. Once dimensionality is controlled for, the metrics agree.

\textbf{Dissociation 2: Linear methods miss the causal signal.} A structured variational autoencoder \citep{narayanaswamy2017learning} discovers choice subspaces that are 3.4$\times$ more causally effective than LDA subspaces, winning in all 73 brain regions. Critically, we validate this against optogenetic silencing data \citep{zatkahaas2021spatiotemporal}: LDA subspace quality is \emph{anti-correlated} with ground-truth causal importance ($\rho = -0.73$), meaning linear methods systematically identify the wrong regions as causally important. The VAE reverses this pattern.

\textbf{An honest caveat.} The interchange intervention accuracy (IIA) metric used to evaluate subspace quality is potentially vacuous for nonlinear methods: both an unconstrained MLP and our structured VAE achieve high IIA on random Gaussian noise (mean $\approx 0.70$). Only linear DAS properly rejects noise (IIA $= 0.16$). This motivates our emphasis on external validation---optogenetic silencing ground truth, multi-task generalization, and engagement orthogonality---over metric-internal IIA comparisons.

These findings converge on a single message: \emph{linear methods fail where they matter most}. The brain regions with the highest effective dimensionality, which are precisely the regions where causal choice signals are strongest, are the regions where LDA, CKA, and linear readouts are least informative. We interpret these results through the neural manifold framework \citep{gallego2017neural}: the choice subspace is the \emph{potent space} for decision output \citep{kaufman2014cortical}, and high-dimensional regions harbor large null spaces that confound linear methods while preserving manifold-level causal structure.

\section{Related Work}

\textbf{Neural manifolds and population-level computation.} \citet{gallego2017neural} argued that the fundamental unit of motor cortical computation is not the single neuron but the \emph{neural mode}---a population-level activation pattern spanning a low-dimensional manifold. Individual corticomotoneuronal cells need not encode any specific movement covariate; what matters is the trajectory through the space of modes. \citet{kaufman2014cortical} formalized this as the potent/null space distinction: preparatory activity lives in a null space orthogonal to the muscle-driving output dimensions, while execution activity enters the potent space. \citet{sadtler2014neural} demonstrated that subjects easily adapt to perturbations within the existing neural manifold but struggle with perturbations that require new modes, confirming that the manifold structure constrains behavior. Our evidence-choice subspace analysis directly tests this framework: the choice subspace is the potent space for decision output, and high-dimensional regions may harbor large null spaces that contain variance but no causal signal.

\textbf{Representational similarity.} CKA \citep{kornblith2019similarity} and RSA \citep{kriegeskorte2008representational} are the standard tools for comparing neural representations. CKA operates on kernel matrices and is therefore linear in the population activity; RSA abstracts from neuron identity by comparing representational dissimilarity matrices (RDMs), preserving relational structure across populations of different size \citep{kriegeskorte2008representational}. \citet{williams2021generalized} showed that many similarity measures are special cases of a generalized shape metric, but all within the linear/kernel family. Our Grassmannian analysis is complementary: RSA compares the global representational structure across all stimuli, while Grassmannian distance compares the subspace structure of the choice-relevant component specifically.

\textbf{Nonlinear methods.} UMAP \citep{mcinnes2018umap} and t-SNE \citep{vandermaaten2008visualizing} are widely used for visualization but less commonly for quantitative comparison. Procrustes analysis on UMAP embeddings has been used for cross-species comparison \citep{deitch2021representational}.

\textbf{Topological and dynamical methods.} Persistent homology has been applied to detect ring-like structures in head direction cells \citep{gardner2022toroidal} and grid cells \citep{rybakken2019decoding}. jPCA \citep{churchland2012neural} extracts rotational dynamics; DSA \citep{ostrow2023beyond} compares dynamical systems.

\textbf{Dimensionality and the cortical hierarchy.} Effective dimensionality scales unboundedly with neuron number \citep{stringer2024unbounded} and changes systematically along the cortical hierarchy \citep{huntenburg2024categorical}. \citet{harvey2024representational} proved that CKA measures average alignment between optimal linear readouts, with sensitivity depending on participation ratio---predicting our empirical finding that high-dimensional regions show low CKA even when manifold shape is conserved.

\textbf{Multimodal parcellation and convergent validity.} \citet{glasser2016multi} defined 180 cortical areas per hemisphere using converging evidence from architecture, function, connectivity, and topography. This multimodal convergence is the neuroscience gold standard for validating that region boundaries are biologically real. Our geometric clustering partially recovers anatomical organization (ARI $= 0.06$, NMI $= 0.30$), suggesting that our metrics capture functional rather than architectonic boundaries.

\textbf{Task transfer structure.} \citet{zamir2018taskonomy} built a computational taxonomy of visual tasks, showing that perceptual and geometric tasks cluster by computational type and that transfer structure is asymmetric. Our cross-region Grassmannian distance matrix is a neural analogue: regions whose choice subspaces are close on the Grassmannian have learned geometrically compatible representations, and the structure of this distance matrix may reveal a functional hierarchy from sensory to association to motor regions.

\textbf{Identifiability and disentanglement.} The semi-supervised disentanglement framework \citep{narayanaswamy2017learning} separates labeled task variables from unlabeled variance. The mechanistic independence condition \citep{sid2026mechanistic} provides identifiability guarantees when the generator Jacobian is sparser in factor-aligned coordinates.

\textbf{IIA vacuity for nonlinear methods.} \citet{sutter2025iia} proved that unconstrained nonlinear alignment makes interchange intervention analysis (IIA) vacuous---any sufficiently flexible nonlinear map can achieve perfect IIA on arbitrary data. The linearity constraint of DAS \citep{geiger2024causal} is what makes IIA non-vacuous. Our structured VAE, despite having a structured latent split, retains enough encoder expressiveness to fit random noise, confirming Sutter et al.'s warning and motivating external validation.

\textbf{Dimension-level causal perturbation.} \citet{pereira2025residual} perturbed coding versus residual subspace dimensions in ALM, finding that residual dimensions causally drive choice through transient amplification. IIA \citep{geiger2024finding} and DAS \citep{geiger2024causal} perform formal causal interventions on learned subspaces in language models.

\textbf{What is missing.} No prior work systematically compares metric families on the same neural data, validates subspace methods against optogenetic ground truth, tests whether the IIA vacuity concern applies to biological subspace estimation, or decomposes the potent/null space variance structure across dozens of brain regions simultaneously.

\section{Methods}

\subsection{Dataset}

We use the Steinmetz et al.\ (2019) dataset: Neuropixels recordings from 10 mice performing a two-alternative visual decision task. The dataset covers 73 brain regions with $\geq 15$ neurons per region (50 with $\geq 2$ recording sessions). We extract trial-averaged activity in a post-stimulus window (150--350ms) and compute binary choice labels (left vs.\ right). All cross-session pairs for the same region constitute the comparison set: 1,316 pairs across 50 multi-session regions.

\subsection{Geometric Invariants}

\textbf{Linear similarity (CKA).} Centered kernel alignment \citep{kornblith2019similarity} on trial $\times$ neuron activity matrices, using the linear kernel.

\textbf{Nonlinear manifold similarity (UMAP Procrustes).} 5-dimensional UMAP embedding \citep{mcinnes2018umap} with Procrustes alignment between sessions.

\textbf{Subspace distance (Grassmannian).} $k$-dimensional LDA subspaces with geodesic distance on $\operatorname{Gr}(k, n)$.

\textbf{Topological invariants.} Vietoris-Rips persistence diagrams via Ripser \citep{bauer2021ripser}.

\textbf{Dynamical invariants.} Simplified jPCA on choice-difference trajectories.

\textbf{Power-law exponent ($\alpha$).} Fit to PCA eigenvalue spectrum in log-log space (ranks 10--50), characterizing dimensionality structure.

\subsection{Structured VAE for Subspace Discovery}

We train a structured variational autoencoder \citep{narayanaswamy2017learning} per region. The encoder maps population activity to two latent groups: $\mathbf{z}_{\text{choice}}$ ($k$ dimensions, $k \in \{1,2,3\}$) trained with an auxiliary classification loss to predict left/right choice, and $\mathbf{z}_{\text{other}}$ ($m = 15$ dimensions) capturing remaining variance. Both groups jointly reconstruct the population activity via the ELBO. Training: 300 epochs, $\alpha_{\text{task}} = 10$, $\beta_{\text{KL}} = 1$, across all 39 sessions.

The encoder weight matrix for $\mathbf{z}_{\text{choice}}$ directly provides the choice subspace directions, replacing the LDA+PCA pipeline.

\subsection{Interchange Intervention Analysis (IIA)}

For each region, we swap the choice-subspace projection between opposite-evidence trial pairs and measure how often a choice classifier flips its prediction. IIA $= 1$ means the subspace is perfectly causal; IIA $= 0.5$ means random. We compute IIA for both the VAE and LDA subspaces, and against a random-direction null.

\subsection{Optogenetic Silencing Validation}

We validate subspace quality against ground-truth causal importance from the Zatka-Haas et al.\ (2021) optogenetic inactivation dataset: 52 cortical coordinates, 47,002 laser trials across 10 mice. For each silencing coordinate mapped to a Steinmetz brain region ($n = 12$ direct matches, $n = 16$ with hierarchical Allen CCF grouping), we compute the behavioral effect size (Euclidean distance in choice-probability space between laser-on and laser-off trials on matched-contrast conditions). The Spearman correlation between silencing effect and subspace IIA provides external validation independent of any metric-internal evaluation.

\subsection{Multi-Task Generalization}

Beyond the primary choice variable, we test subspace discovery on three additional task variables: evidence strength (high vs.\ low contrast), feedback (correct vs.\ error), and stimulus side (left vs.\ right). For each variable, we train a separate structured VAE and compute IIA against the corresponding LDA baseline.

\subsection{Engagement Orthogonality}

We derive an engagement variable from behavioral signatures: disengaged trials (NoGo responses on easy stimuli where the correct response is obvious) vs.\ engaged trials (correct responses on matched stimuli). We estimate engagement and choice subspaces in a pre-stimulus window (0--150ms) and post-stimulus window (250--450ms) respectively, then compute their Grassmannian distance. Near-orthogonality ($d_G \approx \pi/2$ per dimension) would confirm that the VAE disentangles genuinely distinct cognitive variables rather than capturing a single dominant signal.

\section{Results}

\subsection{Dissociation 1: Linear and Nonlinear Metrics Anti-Correlate}

Across 1,316 cross-session pairs and 50 brain regions, CKA and UMAP Procrustes are strongly anti-correlated: Spearman $\rho = -0.85$, 95\% CI $[-0.87, -0.83]$. Motor cortex (MOs) exemplifies the extreme: CKA $= 0.07$ (near-zero linear similarity) but UMAP Procrustes $= 0.98$ (near-maximal nonlinear similarity).

\textbf{The dissociation is robust.} Bootstrap 95\% CI: $[-0.867, -0.830]$. Leave-one-mouse-out jackknife: $\rho \in [-0.872, -0.835]$. Leave-one-region-out: $\rho \in [-0.863, -0.841]$. Neuron sub-sampling: $\rho = -0.81$ ($p < 10^{-7}$). Spectral denoising \citep{spectral2025neurips} \emph{strengthens} the result: $\rho = -0.889$.

\textbf{Dimensionality is the mechanism.} Partial correlation controlling for power-law exponent $\alpha$ reverses the sign: partial $r = +0.44$. Controlling for neuron count alone gives partial $r = +0.13$. Once dimensionality is accounted for, CKA and Procrustes \emph{agree}. This empirically validates the theoretical prediction of \citet{harvey2024representational}.

\textbf{The anti-correlation is metric-specific.} Sliced Wasserstein distance is uncorrelated with dimensionality ($\rho = 0.02$, $p = 0.47$), while SPD Riemannian distance tracks it ($\rho = 0.44$). The dissociation is not a generic ``all metrics disagree'' phenomenon but a specific interaction between kernel-based and manifold-based methods mediated by spectral structure.

\subsection{Spectral Universality and Dynamic Dimensionality}

Singular value spectra are highly conserved across animals (mean $r = 0.94 \pm 0.06$), but temporal modes are idiosyncratic (mean cosine $= 0.09$). The dimensionality \emph{structure} is universal; the \emph{content} is animal-specific.

Dimensionality is not fixed. Under hard task conditions (low contrast), $\alpha$ increases by $+12.4$ (Wilcoxon $p = 3.0 \times 10^{-7}$) while manifold shape is preserved (Procrustes $= 0.90$). Dimensionality compression is an active computational strategy, not a fixed anatomical property.

Rotational dynamics appear in 72/73 regions above a temporal-shuffle null, but strength varies 50$\times$ (range $0.008$--$0.40$). Persistent homology reveals almost no topological structure: only 9/73 regions have $\beta_1 > 0$.

\subsection{Evidence-Choice Subspace Alignment}

For each region, we estimate evidence and choice subspaces and measure their Grassmannian distance. Two strategies emerge:

\begin{enumerate}
\item \textbf{Low-dimensional regions} maintain stable evidence-choice alignment across conditions (high overlap, small choice subspace shift under evidence change).
\item \textbf{High-dimensional regions} show large choice subspace rotations when evidence is experimentally changed ($\rho = -0.59$, permutation $p < 10^{-6}$), maintaining separate, flexibly rotatable representations.
\end{enumerate}

Cross-condition (easy vs.\ hard) subspace rotation also anti-correlates with $\alpha$ ($\rho = -0.65$, 95\% CI $[-0.76, -0.49]$, permutation $p < 0.002$).

\subsection{Dissociation 2: Learned Subspaces Amplify Causal Signal}

\textbf{VAE subspaces are 3.4$\times$ more causal than LDA.} The VAE produces mean IIA $= 0.56$ (SD $0.08$) vs.\ LDA mean IIA $= 0.17$ (SD $0.09$), with Cohen's $d = 4.5$. The VAE wins in 73/73 regions (Wilcoxon $W = 2701$, $p = 5.7 \times 10^{-14}$). It exceeds the random-direction null in all 73 regions (mean null $= 0.23$).

\textbf{Choice subspace dimensionality ablation.} The causal signal is effectively one-dimensional: even $k = 1$ achieves a 3.7$\times$ improvement over LDA (72/73 wins), and performance is stable across $k \in \{1, 2, 3\}$.

\begin{table}[h]
\centering
\small
\caption{Effect of choice subspace dimensionality $k$ on structured VAE IIA across 73 brain regions.}
\label{tab:vae_k_ablation}
\begin{tabular}{lccccc}
\toprule
$k$ & VAE IIA & LDA IIA & Ratio & VAE wins & Choice acc. \\
\midrule
1 & $0.548 \pm 0.093$ & $0.150 \pm 0.091$ & 3.7$\times$ & 72/73 & 0.99 \\
2 & $0.555 \pm 0.088$ & $0.159 \pm 0.097$ & 3.5$\times$ & 73/73 & 0.98 \\
3 & $0.561 \pm 0.080$ & $0.166 \pm 0.094$ & 3.4$\times$ & 73/73 & 0.99 \\
\bottomrule
\end{tabular}
\end{table}

\textbf{Low-dimensional regions benefit most.} Low-$\alpha$ (high-dimensional) regions show a larger VAE improvement ($\Delta$IIA $= 0.46$) than high-$\alpha$ regions ($\Delta$IIA $= 0.33$). VAE IIA anti-correlates with $\alpha$ ($\rho = -0.48$, $p = 1.5 \times 10^{-5}$), reversing the pattern seen with LDA.

\textbf{VAE finds genuinely different subspaces.} The VAE and LDA subspaces are nearly orthogonal: mean Grassmannian distance $= 2.34$ (SD $0.15$), far exceeding $\pi/2 \approx 1.57$ per dimension. The VAE does not refine the LDA solution---it finds a fundamentally different set of directions.

\textbf{Model confidence tracks causal effectiveness.} Posterior uncertainty anti-correlates with IIA ($\rho = -0.26$, $p = 0.028$): regions where the model is most confident are those where interventions are most effective.

\subsection{The IIA Vacuity Problem}
\label{sec:vacuity}

Before interpreting the 3.4$\times$ improvement as evidence that the VAE finds better causal subspaces, we must address a fundamental concern raised by \citet{sutter2025iia}: is IIA a valid metric for nonlinear methods?

We test three methods on real neural data versus random Gaussian noise with matched labels:

\begin{table}[h]
\centering
\small
\caption{IIA on real neural data vs.\ random Gaussian noise (73 regions). A non-vacuous method should show high IIA on real data and chance IIA ($\approx 0.5$) on noise.}
\label{tab:sutter}
\begin{tabular}{lccl}
\toprule
Method & Real IIA & Random IIA & Interpretation \\
\midrule
Unconstrained MLP & --- & $0.70 \pm 0.11$ & Vacuous (fits noise) \\
Structured VAE & $0.69 \pm 0.11$ & $0.70 \pm 0.11$ & Also vacuous on noise \\
Linear DAS & --- & $0.16 \pm 0.09$ & Non-vacuous (rejects noise) \\
\bottomrule
\end{tabular}
\end{table}

The structured VAE achieves comparable IIA on random noise ($0.70$) as on real data ($0.69$). The structured latent split alone does not prevent the encoder from fitting arbitrary labels. This means \emph{the 3.4$\times$ IIA improvement over LDA could in principle reflect encoder expressiveness rather than genuine causal subspace discovery.}

This result motivates two responses: (1) external validation that does not depend on IIA, and (2) ablation experiments that isolate the source of the improvement.

\pending{Exp62: Shuffled-label control. Train VAE with randomly shuffled choice labels on real neural data. If IIA drops to chance, the VAE is learning real structure despite also being able to fit noise. If IIA stays high, the result is vacuous.}

\pending{Exp63: Linear VAE ablation. Compare nonlinear structured VAE, linear structured VAE (same latent split, no hidden-layer nonlinearities), and nonlinear unstructured VAE (single latent, no classification loss). This isolates whether the benefit comes from the structured split or from nonlinearity.}

\subsection{External Validation: Optogenetic Silencing Ground Truth}

The strongest evidence for the VAE's superiority comes not from IIA comparisons but from external validation against optogenetic silencing data---a test that is completely independent of the IIA metric.

Across 12 brain regions matched between Steinmetz and Zatka-Haas et al.\ (2021):

\begin{table}[h]
\centering
\small
\caption{Spearman correlation between subspace quality and optogenetic silencing effect (ground-truth causal importance). Bootstrap BCa CIs from 10,000 resamples.}
\label{tab:silencing}
\begin{tabular}{lccc}
\toprule
Metric & $\rho$ vs.\ silencing & $p$ & 95\% CI \\
\midrule
LDA IIA & $-0.73$ & $0.01$ & $[-0.93, -0.27]$ \\
VAE IIA & $+0.33$ & $0.30$ & --- \\
$\Delta\rho$ (VAE $-$ LDA) & $+1.06$ & $< 0.005$ & $[0.31, 1.51]$ \\
\bottomrule
\end{tabular}
\end{table}

\textbf{LDA subspaces are anti-correlated with causal importance} ($\rho = -0.73$, permutation $p = 0.01$). Regions where optogenetic silencing most disrupts behavior---the ground-truth causally important regions---are the regions where LDA-based IIA is \emph{lowest}. Linear methods systematically identify the wrong regions as causally important.

\textbf{The VAE reverses this pattern.} The VAE correlation is positive ($\rho = +0.33$), though not individually significant at $n = 12$. The critical test is the \emph{difference}: $\Delta\rho = +1.06$, 95\% bootstrap CI $[0.31, 1.51]$, permutation $p < 0.005$. The CI excludes zero, confirming that the VAE provides a significantly better match to optogenetic ground truth than LDA.

\textbf{Hierarchical region expansion.} Using Allen CCF ontology to group Steinmetz sub-regions under parent silencing areas expands the matched set from 12 to 16 regions. The expanded analysis shows VAE $\rho = +0.42$ (perm $p = 0.11$, trending), LDA $\rho = -0.29$, and the delta-rho CI still excludes zero ($\Delta\rho = 0.87$, 95\% CI $[0.09, 1.43]$).

\pending{Exp66: Per-mouse optogenetic validation. Compute per-mouse silencing effects and correlate each mouse's profile with geometry metrics. If the LDA anti-correlation holds per individual mouse (not just pooled), this rules out the concern that outlier mice drive the result.}

\subsection{Multi-Task Generalization}

The VAE advantage is not specific to the choice variable. Across three additional task variables, the VAE wins 73/73 regions for every variable:

\begin{table}[h]
\centering
\small
\caption{VAE vs.\ LDA IIA across four task variables. All 73 regions, all 39 sessions.}
\label{tab:multitask}
\begin{tabular}{lcccc}
\toprule
Task variable & VAE IIA & LDA IIA & VAE wins & Wilcoxon $p$ \\
\midrule
Choice & $0.56 \pm 0.08$ & $0.17 \pm 0.09$ & 73/73 & $5.7 \times 10^{-14}$ \\
Evidence strength & $0.64 \pm 0.13$ & $0.26 \pm 0.10$ & 73/73 & $5.7 \times 10^{-14}$ \\
Feedback & $0.58 \pm 0.11$ & $0.21 \pm 0.09$ & 73/73 & $5.7 \times 10^{-14}$ \\
Stimulus side & $0.66 \pm 0.11$ & $0.29 \pm 0.10$ & 73/73 & $5.7 \times 10^{-14}$ \\
\bottomrule
\end{tabular}
\end{table}

The improvement is largest for evidence strength (2.4$\times$) and smallest for stimulus side (2.3$\times$), but all four show overwhelming VAE dominance. The anti-correlation between $\alpha$ and VAE IIA holds for all task variables: evidence strength ($\rho = -0.75$, $p < 10^{-13}$), feedback ($\rho = -0.70$, $p < 10^{-11}$), stimulus side ($\rho = -0.76$, $p < 10^{-14}$).

\subsection{Engagement-Choice Orthogonality}

The VAE identifies engagement and choice subspaces that are near-orthogonal across all 59 regions with sufficient data:

\begin{itemize}
\item \textbf{Grassmannian distance}: $d_G = 1.82$ (LDA), $d_G = 2.41$ (VAE). Both exceed $\pi/2 \approx 1.57$, confirming orthogonality.
\item \textbf{VAE wins 59/59 regions} for engagement-choice separation.
\item \textbf{Cross-variable IIA}: swapping the engagement subspace between opposite-choice trials produces mean IIA $= 0.17$---near-random, confirming that engagement interventions do not affect choice and vice versa.
\end{itemize}

This demonstrates that the VAE is not merely finding a single dominant nonlinear direction that correlates with everything. It discovers geometrically orthogonal subspaces for genuinely distinct cognitive variables.

\subsection{Cross-Task Subspace Geometry}

We compute pairwise Grassmannian distances between all four task-variable subspaces per region. Two patterns emerge:

\textbf{Choice and stimulus side share structure.} LDA Grassmannian distance between choice and stimulus side is $d_G = 0.585$ (the closest pair), consistent with stimulus side directly driving choice. All other task pairs are near-orthogonal ($d_G \approx 1.3$).

\textbf{VAE separates all task variables more strongly.} The VAE produces larger pairwise distances than LDA in 73/73 regions ($p < 10^{-14}$), with mean $d_G = 2.36$ (VAE) vs.\ $1.20$ (LDA). The VAE finds subspaces that are more orthogonal---better disentangled---than LDA's projections.

\textbf{Multiplexing varies by region.} The most multiplexed regions (lowest mean pairwise distance, most overlapping task subspaces) include OT, LH, and VISa---regions at the interface of multiple processing streams. The least multiplexed (most orthogonal) include PIR, EPd, and AUD---unimodal sensory regions.

\pending{Exp65: Temporal sliding window IIA. Compute IIA in 50ms sliding windows from stimulus onset through response. If the choice subspace IIA rises before the behavioral response, this is evidence of a causal decision signal, not a post-hoc motor artifact.}

\section{Discussion}

\subsection{Linear Methods Fail Where They Matter Most}

Our central claim is not merely that nonlinear methods outperform linear ones---this has been shown before in various settings. The novel finding is the \emph{interaction} with dimensionality: linear methods fail most severely in the highest-dimensional brain regions, and these are precisely the regions where optogenetic silencing reveals the strongest causal role in behavior.

This creates a systematic blind spot. A researcher using LDA to identify causally important regions would conclude that low-dimensional sensory regions carry the strongest choice signals (because LDA IIA is highest there), while high-dimensional association and hippocampal regions are causally weak. The optogenetic ground truth shows the opposite: the high-dimensional regions are the most causally important, but their causal subspaces are invisible to linear methods.

\subsection{The Neural Manifold Hypothesis as Unifying Framework}

Our results are naturally interpreted through the neural manifold framework of \citet{gallego2017neural}. The evidence subspace identified by the VAE corresponds to the \emph{potent space} for choice output \citep{kaufman2014cortical}---the subspace of population activity that causally drives the downstream binary decision. The remaining variance, which may be high-dimensional in association regions like MOs, lives in the \emph{null space}---it is present, it contributes to CKA dissimilarity, but it does not drive behavior.

This potent/null decomposition explains the CKA--Procrustes dissociation. A region like MOs shows CKA $\approx 0$ because the high-dimensional null space dominates the kernel matrix, making sessions appear dissimilar. But UMAP Procrustes $\approx 1$ because the manifold \emph{shape}---including the low-dimensional choice-potent subspace embedded within it---is conserved. CKA tests whether two regions share the same \emph{task-specific coordinates} (same directions of variance); Procrustes tests whether they share the same \emph{universal manifold structure} (same geometric shape). High Procrustes with low CKA means: same universal manifold, different subspace sampled within it.

The within-manifold / outside-manifold learning distinction of \citet{sadtler2014neural} makes a testable prediction: if the VAE's choice subspace represents a within-manifold representation, then IIA interventions should be robust to noise (the manifold provides a scaffold for downstream interpretation). Conversely, if choice encoding is outside-manifold for some regions, IIA should be noisier. This predicts a correlation between manifold dimensionality and IIA reliability across regions.

\pending{Exp67: Potent/null space decomposition. For each region, project activity onto the choice subspace (potent) and its complement (null). Compute variance fractions. Test whether high-dimensional regions have larger null spaces and whether null-space variance is uncorrelated with choice.}

\subsection{The IIA Vacuity Problem and External Validation}

We take the Sutter dilemma seriously. Our structured VAE achieves IIA $= 0.70$ on random Gaussian noise---the same as an unconstrained MLP. This means IIA alone cannot distinguish genuine causal subspace discovery from encoder flexibility. The structured latent split (z\_choice/z\_other) does not provide sufficient constraint to prevent overfitting.

Three lines of evidence nevertheless support the VAE's validity:

\begin{enumerate}
\item \textbf{Optogenetic validation} ($\S$5.6): The VAE-silencing correlation is positive while LDA's is strongly negative, and the difference is significant. This test is completely independent of IIA.
\item \textbf{Cross-task disentanglement} ($\S$5.8): The VAE discovers near-orthogonal subspaces for engagement vs.\ choice, with cross-variable IIA at chance ($0.17$). A vacuous method would not produce orthogonality between genuinely distinct cognitive variables.
\item \textbf{Multi-task generalization} ($\S$5.7): The VAE advantage holds for all four task variables with the same dimensionality interaction pattern. A single overfitting mechanism would not generalize this consistently.
\end{enumerate}

\pending{Exp62 (shuffled-label control) will provide the definitive test: if the VAE's IIA drops when trained on randomly shuffled labels, the encoder IS learning real structure despite also being able to fit noise on random data.}

\pending{Exp63 (linear VAE ablation) will isolate whether the benefit comes from the structured latent split or from encoder nonlinearity.}

\subsection{Dimensionality as Computational Strategy}

The task-difficulty modulation ($\Delta\alpha = +12.4$, $p = 3 \times 10^{-7}$) while preserving manifold shape demonstrates that dimensionality is an active computational variable, not a fixed property. Under difficult conditions, circuits compress into fewer dimensions---concentrating variance along task-relevant axes. The evidence-choice subspace analysis reveals two strategies: low-dimensional regions route through shared subspaces (efficient, inflexible), while high-dimensional regions maintain separate, rotatable representations (flexible, requiring nonlinear readout).

\subsection{Toward Neural Taskonomy and RSA of Subspaces}

Our pairwise Grassmannian distances between regions' choice subspaces define a \emph{Grassmannian dissimilarity matrix} (GDM)---an analogue of the representational dissimilarity matrix (RDM) used in RSA \citep{kriegeskorte2008representational}, but operating on subspaces rather than activity patterns. Where RSA asks ``do two conditions evoke similar population responses?'', the GDM asks ``do two regions use similar geometric strategies to encode choice?'' Correlating the GDM with CKA-based and Procrustes-based RDMs, anatomical distance, and functional connectivity would reveal whether subspace similarity captures biologically meaningful organization that standard metrics miss.

This extends naturally to a \emph{neural taskonomy} \citep{zamir2018taskonomy}. Just as Zamir et al.\ discovered that visual tasks form a computational taxonomy with asymmetric transfer, the GDM may reveal a functional hierarchy in which sensory regions' subspaces are close to each other but distant from motor regions', with association regions in between. The structure of this graph---and whether it predicts which region's lesion most disrupts another region's decoding---is a direct test of the manifold framework applied to brain-wide circuit organization.

\pending{Exp68: Subspace dissimilarity matrix. Build a $73 \times 73$ GDM from pairwise Grassmannian distances between regions' choice subspaces. Correlate with CKA RDM, Procrustes RDM, and anatomical distance via Mantel tests. Test whether the GDM recovers functional hierarchy (sensory $\to$ association $\to$ motor).}

\subsection{Practical Implications}

For neural data analysis, our findings imply a simple diagnostic: compute the power-law exponent $\alpha$ before selecting analysis methods. For $\alpha > 3$ (low-dimensional), linear methods (CKA, LDA) are informative. For $\alpha < 1$ (high-dimensional), nonlinear methods are necessary---and linear methods may be actively misleading. For intermediate regimes, both should be reported.

For the interpretability community, the IIA vacuity result is a cautionary finding. IIA is widely used to evaluate causal subspace quality in language models; our results suggest that external validation (behavioral perturbation, out-of-distribution generalization) is necessary alongside metric-internal evaluation, particularly for nonlinear methods.

\subsection{Limitations}

\textbf{Dataset scale.} 10 mice, 73 regions. The IBL Brain-Wide Map (139 mice) would enable per-region significance testing. \pending{Exp46: IBL cross-dataset transfer.}

\textbf{Silencing sample size.} $n = 12$ direct matches ($n = 16$ with hierarchical grouping) for the optogenetic validation. \pending{Exp66 (per-mouse analysis) and exp64 (hierarchical expansion) address this.}

\textbf{IIA vacuity.} The structured VAE's IIA is potentially vacuous (\S\ref{sec:vacuity}). We rely on external validation, but \pending{exp62 (shuffled control) and exp63 (linear ablation) will provide the definitive internal controls.}

\textbf{No temporal resolution.} All analyses use trial-averaged activity (150--350ms window). \pending{Exp65 (temporal IIA) will test whether the choice subspace emerges before the behavioral response.}

\textbf{Task specificity.} All results are for binary visual choice. Other task types (navigation, working memory) may show different patterns.

\section{Conclusion}

We establish two geometric dissociations in neural population codes. First, linear and nonlinear similarity metrics anti-correlate ($\rho = -0.85$), mediated by effective dimensionality. Second, linear subspace methods systematically underestimate causal structure by 3.4$\times$ and are \emph{anti-correlated} with optogenetic ground-truth causal importance ($\rho = -0.73$). These dissociations converge: the brain regions where linear methods fail most severely are precisely the regions where causal signals are strongest. Structured nonlinear methods recover these signals, producing subspaces that generalize across four task variables, disentangle engagement from choice, and better predict optogenetic silencing effects. However, the IIA metric itself is potentially vacuous for nonlinear methods, motivating external validation as the gold standard for subspace quality assessment.

\bibliographystyle{tmlr}
\bibliography{references}

\appendix
\section{Sequential Discovery and Multiple Comparisons}
\label{app:sequential}

\begin{table}[h]
\centering
\small
\begin{tabular}{@{}clcccc@{}}
\toprule
\textbf{Stage} & \textbf{Test} & \textbf{$\rho$} & \textbf{$p$} & \textbf{$N_k$} & \textbf{Survives} \\
\midrule
\multicolumn{6}{l}{\emph{Stage 1: Core metric comparison}} \\
1 & CKA--Procrustes anti-correlation & $-0.85$ & $\approx 0$ & 1 & \checkmark \\
\midrule
\multicolumn{6}{l}{\emph{Stage 2: Mechanism identification}} \\
2 & Dimensionality mediates (partial $r$) & $+0.44$ & $< 10^{-6}$ & 2 & \checkmark \\
\midrule
\multicolumn{6}{l}{\emph{Stage 3: Causal subspace discovery}} \\
3 & VAE $>$ LDA IIA (73/73 regions) & --- & $5.7 \times 10^{-14}$ & 3 & \checkmark \\
4 & Silencing vs.\ LDA IIA & $-0.73$ & $0.01$ & 4 & \checkmark \\
5 & $\Delta\rho$ (VAE $-$ LDA) vs.\ silencing & $+1.06$ & $< 0.005$ & 5 & \checkmark \\
\midrule
\multicolumn{6}{l}{\emph{Stage 4: Generalization}} \\
6 & Multi-task VAE wins (4 variables) & --- & $< 10^{-13}$ & 6 & \checkmark \\
7 & Engagement-choice orthogonality & $d_G = 1.82$ & --- & 7 & \checkmark \\
\midrule
\multicolumn{6}{l}{\emph{Stage 5: IIA vacuity and controls}} \\
8 & VAE random IIA $\approx$ MLP random IIA & $0.70$ & $p = 0.38$ (no diff) & 8 & (diagnostic) \\
\bottomrule
\end{tabular}
\caption{Sequential discovery table. All seven core claims (stages 1--4) survive Holm correction. Stage 5 is a diagnostic test, not a directional hypothesis.}
\label{tab:sequential}
\end{table}

\section{Robustness Controls}
\label{app:robustness}

\begin{table}[h]
\centering
\small
\caption{Robustness of the CKA--Procrustes anti-correlation.}
\label{tab:robustness}
\begin{tabular}{@{}lccl@{}}
\toprule
\textbf{Control} & \textbf{$\rho$} & \textbf{$p$} & \textbf{$n$} \\
\midrule
Raw (full data) & $-0.850$ & $\approx 0$ & 1,316 pairs \\
Bootstrap 95\% CI & $[-0.867, -0.830]$ & --- & 10,000 resamples \\
Leave-one-mouse-out (range) & $[-0.872, -0.835]$ & all sig. & 10 mice \\
Leave-one-region-out (range) & $[-0.863, -0.841]$ & all sig. & 50 regions \\
Neuron sub-sampling & $-0.814$ & $< 10^{-7}$ & 29 pairs \\
Spectral denoising (MP) & $-0.889$ & $\approx 0$ & 1,316 pairs \\
Trial-label shuffling & $-0.395$ & $< 10^{-50}$ & 1,316 pairs \\
\midrule
\multicolumn{4}{l}{\emph{Partial correlations (Pearson $r$)}} \\
Controlling neuron count & $+0.134$ & --- & sign reversal \\
Controlling $\alpha$ & $+0.437$ & --- & sign reversal \\
Controlling both & $+0.196$ & --- & sign reversal \\
\bottomrule
\end{tabular}
\end{table}

\section{Metric Family Comparison}
\label{app:metrics}

\begin{table}[h]
\centering
\small
\caption{Correlation of each metric family with dimensionality ($\alpha$) and with CKA/Procrustes.}
\label{tab:metrics}
\begin{tabular}{@{}lccccc@{}}
\toprule
\textbf{Metric} & \textbf{vs.\ $\alpha$} & \textbf{$p$} & \textbf{vs.\ CKA} & \textbf{vs.\ Procrustes} \\
\midrule
CKA (linear) & $-0.68$ & $< 10^{-4}$ & --- & $-0.85$ \\
UMAP Procrustes & $+0.68$ & $< 10^{-4}$ & $-0.85$ & --- \\
Sliced Wasserstein & $+0.02$ & $0.47$ & $+0.10$ & $-0.09$ \\
SPD Riemannian & $+0.44$ & $< 10^{-4}$ & --- & --- \\
Fisher-Rao & $+0.40$ & $< 10^{-4}$ & --- & --- \\
\bottomrule
\end{tabular}
\end{table}

\section{Evidence-Choice Subspace Alignment}
\label{app:alignment}

\begin{table}[h]
\centering
\small
\caption{Top and bottom 5 regions by evidence-choice subspace alignment.}
\label{tab:routing}
\begin{tabular}{@{}lccc@{}}
\toprule
\textbf{Region} & \textbf{$d_G(V_{\text{ev}}, V_{\text{ch}})$} & \textbf{$\alpha$} & \textbf{Interpretation} \\
\midrule
\multicolumn{4}{l}{\emph{Highest alignment (shared subspace)}} \\
OT & 0.197 & 104.2 & evidence-choice routing \\
IC & 0.409 & 13.1 & sensorimotor integration \\
MEA & 0.453 & 3.0 & evidence-choice routing \\
LH & 0.471 & 3.9 & evidence-choice routing \\
SCm & 0.486 & 15.6 & motor output \\
\midrule
\multicolumn{4}{l}{\emph{Lowest alignment (separated subspaces)}} \\
VAL & 1.163 & 1.7 & thalamic relay \\
LSc & 1.216 & 2.4 & basal ganglia \\
COA & 1.313 & 0.7 & amygdala \\
CL & 1.381 & 27.7 & thalamic relay \\
EPd & 1.397 & 43.6 & amygdala/piriform \\
\bottomrule
\end{tabular}
\end{table}

\section{Topological and Dynamical Invariants}
\label{app:topo_dyn}

\textbf{Persistent homology.} Only 9/73 regions have $\beta_1 > 0$. Grassmannian vs.\ topological Wasserstein $H_1$ distance: $\rho = 0.25$, $p = 0.52$ ($n = 9$).

\textbf{Rotational dynamics.} 72/73 regions show rotation above null. Mean strength $= 0.053 \pm 0.060$; range $0.008$--$0.40$.

\textbf{Spectral universality.} Singular value spectra: mean $r = 0.94 \pm 0.06$ (1,316 pairs). Temporal modes: mean cosine $= 0.09 \pm 0.07$.

\section{Tucker Decomposition}
\label{app:tucker}

Tucker decomposition on (trials $\times$ neurons $\times$ time) tensors: time factors are most conserved (mean similarity $0.174$), followed by trial ($0.106$) and neuron ($0.067$). Tucker $R^2$ correlates modestly with $\alpha$ ($\rho = 0.275$).

\end{document}
